% chapter: 付録（問いの連鎖）

\chapter{解法の背景と広がり}\label{章：解法の背景と広がり}
\begin{bookcolumns}

  付録では,本編で選んだ量や変形を別の表現から見直す.
  \begin{itemize}
    \item 連立型・隣接3項間の更新をまとめるなら,\bookterm{matrix}{行列}の基本へ.
    \item \bookterm{difference}{階差}・総和と微積分をつなぐなら,\sectionref*{節：和と差から微積分へ}へ.微積分を使う部分は任意に読める.
    \item \bookterm{sequence}{数列}全体の係数を扱うなら,\sectionref*{節：数列をまとめて見る}へ.\bookterm{generating}{母関数}は級数の収束を仮定せずに扱う.
  \end{itemize}
  最初の一巡では,本編の操作との対応と,新しく見えた性質を確かめればよい.

  % BOOK-FILE-ID: f-de418d1ee182
% 付録：行列の基本

\section{行列の基本}\label{節：行列の基本}

\sectionref*{節：連立型}では,二つの\bookterm{sequence}{数列}から和と差を作り,漸化式を簡単にした.
例えば
\[
  x_{n+1}=2x_n+y_n,\qquad y_{n+1}=x_n+2y_n
\]
なら,\,$(x_1,y_1)=(1,0)$から$(x_2,y_2)=(2,1)$,次に$(x_3,y_3)=(5,4)$と進む.
毎回の計算には2つの値が両方必要なので,いっそ2つの値を1組として記録し,その組をまとめて計算してしまおう.
そのとき,漸化式の係数をまとめて書くための道具が\bookterm{matrix}{行列}である.
まずは行列の掛け方を学び,そこから連立型で行った「等比型への分解」が何をしていたのかを見ていく.

\subsection{行列とベクトル}

\bookterm{matrix}{行列}とは,数を長方形に並べ,全体を括弧で囲んだものである.
たとえば
\[
  A=\begin{pmatrix}2&1\\1&2\end{pmatrix}
\]
は二行二列の\bookterm{matrix}{行列}である.横の並びを\textbf{行},縦の並びを\textbf{列},各数を\textbf{成分}という.
行は上から,列は左から数え,この例の第一行・第二列の成分は$1$である.
行列の定数倍は各成分をその定数倍する.

また,二つの数を縦に並べた
\[
  \boldsymbol{v}=\begin{pmatrix}x\\y\end{pmatrix}
\]
を\textbf{\bookterm{column-vector}{列ベクトル}}と呼ぶ.
太字の$\boldsymbol{v}$は二つの値をまとめた記号で,平面上の点$(x,y)$への矢印とも読める.
0でない二つのベクトルが直交するとは,矢印が垂直に交わることである.
成分が$(u_1,u_2)$と$(v_1,v_2)$なら,直交する条件は$u_1v_1+u_2v_2=0$である.

\bookterm{column-vector}{列ベクトル}の足し算と定数倍は,成分ごとに行う.
\[
  \begin{pmatrix}x\\y\end{pmatrix}
  +\begin{pmatrix}u\\v\end{pmatrix}
  =\begin{pmatrix}x+u\\y+v\end{pmatrix},
  \qquad
  c\begin{pmatrix}x\\y\end{pmatrix}
  =\begin{pmatrix}cx\\cy\end{pmatrix}.
\]
たとえば$\begin{pmatrix}3\\1\end{pmatrix}+\begin{pmatrix}1\\-1\end{pmatrix}=\begin{pmatrix}4\\0\end{pmatrix}$である.
等号も成分ごとの等しさを表す.
一般の二成分のベクトルも
\[
  \begin{pmatrix}x\\y\end{pmatrix}
  =\frac{x+y}{2}\begin{pmatrix}1\\1\end{pmatrix}
  +\frac{x-y}{2}\begin{pmatrix}1\\-1\end{pmatrix}
\]
と分けられる.本編の和と差は,この二つの係数を求める操作に対応する.
\bookproblemref{practice-basic}{7}{標準問7}でも使うこの分解が,後で固有値に結び付く.

\subsubsection{行列を掛ける}

先ほどの$A$を\bookterm{column-vector}{列ベクトル}に掛けると
\[
  \begin{pmatrix}2&1\\1&2\end{pmatrix}
  \begin{pmatrix}x\\y\end{pmatrix}
  =
  \begin{pmatrix}2x+y\\x+2y\end{pmatrix}.
\]
となる.各行の数を$x,y$にそれぞれ掛けて足す.
たとえば
\[
  A\begin{pmatrix}1\\0\end{pmatrix}
  =\begin{pmatrix}2\\1\end{pmatrix},
  \qquad
  A\begin{pmatrix}2\\1\end{pmatrix}
  =\begin{pmatrix}5\\4\end{pmatrix}.
\]
$A$は$x,y$を$2x+y,x+2y$へ移す規則を表す.
\BookFigMatrixMap

一般には
\[
  \begin{pmatrix}p&q\\r&s\end{pmatrix}
  \begin{pmatrix}x\\y\end{pmatrix}
  =\begin{pmatrix}px+qy\\rx+sy\end{pmatrix}
\]
と定める.連立式$x'=px+qy,\ y'=rx+sy$を積一つにまとめたもので,\,$x',y'$は$A$を掛けた後の値を表す.
このように行列を掛けてベクトルを移す操作を一次変換という.

また,成分ごとに展開すれば
\[
  A(c\boldsymbol{u}+d\boldsymbol{v})
  =cA\boldsymbol{u}+dA\boldsymbol{v}
\]
が成り立つ.
つまり,足してから$A$を掛けても,別々に$A$を掛けてから足しても同じになる.後で初期値を分ける際に使う.

\subsection{行列の積}

$A$を2回掛ける操作をまとめた\bookterm{matrix}{行列}を$A^2$と書く.文字$x,y$で計算すると
\[
  \begin{aligned}
    A\left(A\begin{pmatrix}x\\y\end{pmatrix}\right)
    &=A\begin{pmatrix}2x+y\\x+2y\end{pmatrix}\\
    &=\begin{pmatrix}5x+4y\\4x+5y\end{pmatrix}.
  \end{aligned}
\]
したがって$A^2=\begin{pmatrix}5&4\\4&5\end{pmatrix}$である.各成分の二乗とは異なる.

異なる\bookterm{matrix}{行列}$A,B$でも,まず$B$,次に$A$を掛ける操作を$AB$と書く.
$B$の各列に$A$を掛けた結果を並べると,積$AB$が得られる.
具体的には
\[
  \begin{aligned}
    &\begin{pmatrix}p&q\\r&s\end{pmatrix}
    \begin{pmatrix}u&v\\w&z\end{pmatrix}\\
    &\qquad=
    \begin{pmatrix}
      pu+qw&pv+qz\\
      ru+sw&rv+sz
    \end{pmatrix}.
  \end{aligned}
\]
右の\bookterm{matrix}{行列}から先に作用するため,一般には$AB$と$BA$は一致しない.

何も変えない変換を表す行列は
\[
  I=\begin{pmatrix}1&0\\0&1\end{pmatrix},
  \qquad
  I\begin{pmatrix}x\\y\end{pmatrix}
  =\begin{pmatrix}x\\y\end{pmatrix}
\]
で表される.
$I$を\textbf{\bookterm{identity}{単位行列}}という.
数の累乗で$a^0=1$とするのに対応して,\bookterm{matrix}{行列}では$A^0=I$と定める.

\subsection{回転と複素数の掛け算}\label{節：回転と複素数の掛け算}

同じ\bookterm{matrix}{行列}を繰り返し掛けると何が起こるのか,図で見やすい例を1つ用意しておこう.
原点を中心に,点を角$\theta$だけ回す操作である.

長さ1の点$(\cos\varphi,\sin\varphi)$を角$\theta$だけ回すと,\,$(\cos(\varphi+\theta),\sin(\varphi+\theta))$に移る.
\bookterm{addition}{加法定理}で展開すると
\[
  \begin{aligned}
    \cos(\varphi+\theta)&=\cos\theta\cos\varphi-\sin\theta\sin\varphi,\\
    \sin(\varphi+\theta)&=\sin\theta\cos\varphi+\cos\theta\sin\varphi
  \end{aligned}
\]
である.\,$\cos\varphi,\sin\varphi$に掛かる係数を並べると
\[
  R(\theta)=\begin{pmatrix}\cos\theta&-\sin\theta\\\sin\theta&\cos\theta\end{pmatrix}
\]
となり,\,$R(\theta)\begin{pmatrix}\cos\varphi\\\sin\varphi\end{pmatrix}=\begin{pmatrix}\cos(\varphi+\theta)\\\sin(\varphi+\theta)\end{pmatrix}$と書ける.
長さ$r$の点$r(\cos\varphi,\sin\varphi)$も,両辺を$r$倍すれば同じように回る.
この$R(\theta)$を\textbf{回転行列}という.

回転を2回続ければ角$2\theta$の回転になるので,\,$R(\theta)^2=R(2\theta)$である.
繰り返せば
\[
  R(\theta)^n=R(n\theta)
\]
となる.同じ行列を繰り返し掛けることが,一定の角を何度も足すことになるのだ.

この行列は,複素数の掛け算にも現れる.
複素数$a+bi$を$x+yi$に掛けると
\[
  (a+bi)(x+yi)=(ax-by)+(bx+ay)i
\]
となる.実部と虚部を縦に並べたベクトルで見ると
\[
  \begin{pmatrix}ax-by\\bx+ay\end{pmatrix}
  =\begin{pmatrix}a&-b\\b&a\end{pmatrix}\begin{pmatrix}x\\y\end{pmatrix}
\]
である.ここで$a+bi=r(\cos\theta+i\sin\theta)$と極形式で書けば,
\[
  \begin{pmatrix}a&-b\\b&a\end{pmatrix}=rR(\theta)
\]
となる.つまり,複素数を掛けることは,\textbf{角$\theta$の回転と$r$倍の拡大・縮小}を表す行列を掛けることなのである.
\sectionref*{節：特性方程式の解が虚数の場合}で複素数の冪が回転として現れたのは,このためだったのだ.

\subsection{固有ベクトルと固有値}

毎回\bookterm{matrix}{行列}の積を計算するのは手間がかかる.
そこで,掛けても定数倍になるだけの\bookterm{column-vector}{列ベクトル}を探そう.
先ほどの$A$について
\[
  \begin{aligned}
    A\begin{pmatrix}1\\1\end{pmatrix}
    &=\begin{pmatrix}3\\3\end{pmatrix}
    =3\begin{pmatrix}1\\1\end{pmatrix},\\
    A\begin{pmatrix}1\\-1\end{pmatrix}
    &=\begin{pmatrix}1\\-1\end{pmatrix}
  \end{aligned}
\]
である.
一方は毎回$3$倍され,他方は変わらないので
\[
  A^k\begin{pmatrix}1\\1\end{pmatrix}
  =3^k\begin{pmatrix}1\\1\end{pmatrix},
  \qquad
  A^k\begin{pmatrix}1\\-1\end{pmatrix}
  =\begin{pmatrix}1\\-1\end{pmatrix}
\]
となる.

一般に,0でない\bookterm{column-vector}{列ベクトル}$\boldsymbol{u}$が
\[
  A\boldsymbol{u}=\lambda\boldsymbol{u}
\]
を満たすとき,\,$\boldsymbol{u}$を$A$の\textbf{\bookterm{eigenvector}{固有ベクトル}},\,$\lambda$を対応する\textbf{\bookterm{eigenvalue}{固有値}}という.
\bookterm{eigenvalue}{固有値}は,そのベクトルに対して$A$を1回掛けると何倍になるかを表す数である.
固有ベクトルを0でない数で定数倍したものも,同じ固有値の固有ベクトルである.原点を通る直線(方向)だけが決まり,長さや向きは自由に選べる.
零ベクトルはどんな$\lambda$でも式を満たすため除く.負の\bookterm{eigenvalue}{固有値}では向きが反転し,\,$0$なら1回掛けただけで零ベクトルになる.

\subsubsection{固有ベクトルに分ける}

\bookterm{eigenvector}{固有ベクトル}が見つかると,ほかのベクトルへの作用も計算しやすくなる.
たとえば
\[
  \begin{pmatrix}1\\0\end{pmatrix}
  =\frac12\begin{pmatrix}1\\1\end{pmatrix}
  +\frac12\begin{pmatrix}1\\-1\end{pmatrix}
\]
である.
\booknavref{解釈：三項間と解全体}{隣接3項型}で一般項を2つの解の和に分けたように,ここでは$(1,0)$というベクトルを,二つの\bookterm{column-vector}{列ベクトル}の\bookterm{combination}{線形結合}で表している.
2つの部分に別々に$A$を$k$回掛ければ
\[
  \begin{aligned}
    A^k\begin{pmatrix}1\\0\end{pmatrix}
    &=\frac{3^k}{2}\begin{pmatrix}1\\1\end{pmatrix}
    +\frac12\begin{pmatrix}1\\-1\end{pmatrix}\\
    &=\begin{pmatrix}\frac{3^k+1}{2}\\\frac{3^k-1}{2}\end{pmatrix}.
  \end{aligned}
\]
本編の等比型への帰着と同じく,一定の倍数で変わる部分に分ければ,大きな累乗も計算できる.
図では,定数倍される二つのベクトルを$\boldsymbol v,\boldsymbol w$,
一般の方向の例を$\boldsymbol u=(1,0)$として描く.
\BookFigEigen

一般に,二つの異なる\bookterm{eigenvalue}{固有値}$\lambda_1,\lambda_2$に対応する\bookterm{eigenvector}{固有ベクトル}$\boldsymbol{u}_1,\boldsymbol{u}_2$があれば,どんな二成分のベクトルも
\[
  \boldsymbol{v}=c\boldsymbol{u}_1+d\boldsymbol{u}_2
\]
と表せる.
異なる\bookterm{eigenvalue}{固有値}に対応する二つの\bookterm{eigenvector}{固有ベクトル}は,互いの定数倍にならないからである.
係数$c,d$は,上下の成分を比較した連立方程式で求められる.
その後は
\[
  A^k\boldsymbol{v}
  =c\lambda_1^k\boldsymbol{u}_1+d\lambda_2^k\boldsymbol{u}_2
\]
で計算できる.
$k=0$のときは元の分解式を読み,\,$k\ge1$で\bookterm{eigenvalue}{固有値}が$0$なら対応する部分は0になる.

\subsection{行ベクトルと左固有ベクトル}

数を横に並べた$L=(\alpha,\beta)$を\textbf{\bookterm{row-vector}{行ベクトル}}という.
\bookterm{column-vector}{列ベクトル}$\boldsymbol{\ell}$を横に並べ直す操作を\textbf{\bookterm{transpose}{転置}}といい,\,$\boldsymbol{\ell}^{\mathsf T}$と書く.
たとえば$\boldsymbol{\ell}=\begin{pmatrix}\alpha\\\beta\end{pmatrix}$なら,\,$\boldsymbol{\ell}^{\mathsf T}=(\alpha,\beta)$である.
行列の転置も,行と列を入れ替える操作であり,\,$A^{\mathsf T}$と書く.
\[
  \begin{pmatrix}p&q\\r&s\end{pmatrix}^{\mathsf T}
  =\begin{pmatrix}p&r\\q&s\end{pmatrix}.
\]
\bookterm{row-vector}{行ベクトル}と\bookterm{column-vector}{列ベクトル}の積は,対応する成分を掛けて足す.
\[
  (\alpha,\beta)\begin{pmatrix}x\\y\end{pmatrix}
  =\alpha x+\beta y.
\]
\bookterm{row-vector}{行ベクトル}を\bookterm{matrix}{行列}に掛けるときは,\bookterm{matrix}{行列}の各列との積を横に並べる.
たとえば,先ほどの$A$について
\[
  \begin{aligned}
    (1,\,1)A&=(3,\,3)=3(1,\,1),\\
    (1,-1)A&=(1,-1)
  \end{aligned}
\]
である.
一般に,0でない\bookterm{row-vector}{行ベクトル}$L$が$LA=\lambda L$を満たすとき,\,$L$を\textbf{\bookterm{left-eigen}{左固有ベクトル}}という.
この場合,\bookterm{column-vector}{列ベクトル}$\boldsymbol{v}$から作った数$L\boldsymbol{v}$は,\,$A$を掛けた後には
\[
  L(A\boldsymbol{v})=(LA)\boldsymbol{v}
  =\lambda L\boldsymbol{v}
\]
となる.
たとえば$L=(1,\,1)$なら,\,$x+y$という和が$A$を掛けるたびに$3$倍される.
\bookterm{eigenvector}{固有ベクトル}は定数倍になる\bookterm{column-vector}{列ベクトル}を表し,\bookterm{left-eigen}{左固有ベクトル}は定数倍になる量を取り出す係数を表す.

\subsection{行列式と固有値}

\bookterm{eigenvalue}{固有値}の条件をまとめるため,\bookterm{matrix}{行列}式を導入する.
二行二列の場合には
\[
  \det\begin{pmatrix}a&b\\c&d\end{pmatrix}=ad-bc
\]
と定める.
\bookterm{matrix}{行列}そのものではなく,四つの成分から計算した一つの数である.

二つの式$ax+by=0,\ cx+dy=0$について,\,$ad-bc\ne0$なら解は$x=y=0$だけになる.
たとえば,第一の式を$d$倍したものから第二の式を$b$倍したものを引くと,\,$(ad-bc)x=0$となる.
同様に$(ad-bc)y=0$も得られる.
逆に$ad-bc=0$なら,0でない解が存在する.
実際,\,$(a,b)\ne(0,\,0)$なら$(x,y)=(b,-a)$が両方の式を満たす.
$(a,b)=(0,\,0)$の場合も,残る一つの式から0でない解を選べる.

\bookterm{eigenvalue}{固有値}の式$A\boldsymbol{u}=\lambda\boldsymbol{u}$は
\[
  (\lambda I-A)\boldsymbol{u}=\boldsymbol{0}
\]
と書き直せる.
\bookterm{matrix}{行列}の引き算は対応する成分ごとに行うので,\,$A=\begin{pmatrix}p&q\\r&s\end{pmatrix}$なら
\[
  \lambda I-A
  =\begin{pmatrix}\lambda-p&-q\\-r&\lambda-s\end{pmatrix}.
\]
0でない解が存在する条件は
\[
  \begin{aligned}
    \det(\lambda I-A)
    &=(\lambda-p)(\lambda-s)-qr\\
    &=\lambda^2-(p+s)\lambda+(ps-qr)=0.
  \end{aligned}
\]
本書ではこれを\bookterm{matrix}{行列}$A$の\textbf{固有値の方程式}と呼ぶ.\footnote{行列の\bookterm{characteristic}{特性方程式}と呼ぶのが一般的である.本編で「特性方程式」と呼んだ方程式には,この固有値の方程式にあたるものと,不動点の値を求める方程式にあたるものがあるので,付録では呼び分けておく.}
$\lambda$を求めた後,成分の連立方程式から対応する\bookterm{eigenvector}{固有ベクトル}を求められる.
たとえば$A=\begin{pmatrix}2&1\\1&2\end{pmatrix}$では
\[
  \lambda^2-4\lambda+3
  =(\lambda-3)(\lambda-1)=0
\]
だから,\bookterm{eigenvalue}{固有値}は$3,\,1$である.
これは数値例で確かめた倍数と一致する.

\bookterm{recurrence}{漸化式}では,次の項を計算するために記録しておく値の組を\textbf{状態}と呼ぶ.
また,\bookterm{matrix}{行列}式が$0$でないときは,\,$A$を掛けた後の2つの値から掛ける前の2つの値を一意に求め直せる.
このように逆向きに戻せる\bookterm{matrix}{行列}を\textbf{可逆な\bookterm{matrix}{行列}}という.

次節では,本編で引いた定数や作った比・和を,\bookterm{matrix}{行列}の中で見直す.

  % 付録：本編の置き換えを行列で捉え直す

\section{漸化式と行列}\label{節：漸化式と行列}

本編で作った\bookterm{geometric}{等比数列}を,\bookterm{matrix}{行列}で見直す.
\sectionref{節：特性方程式型}の\bookterm{fixed}{不動点}から,一次分数型,隣接3項間型,連立型へ進もう.
以下の「二階」は,次の項に直前の二項が必要なことをいい,本編の隣接3項間型に対応する.

\subsection{一つの値を更新する}\label{節：一つの値を更新する}

本編の例題$a_{n+1}=2a_n+1$では,毎回定数$1$が加わる.
その$1$も,$a_n$と一緒に記録する.
二つ目の成分を常に$1$に保てば,足し算もその成分の定数倍として書ける.
\[
  \begin{pmatrix}2&1\\0&1\end{pmatrix}
  \begin{pmatrix}a_n\\1\end{pmatrix}
  =\begin{pmatrix}2a_n+1\\1\end{pmatrix}
\]
と書ける.下の成分は毎回$1$のままで,上の成分だけが元の規則で更新される.

一般の$a_{n+1}=pa_n+q$ $(p\ne0,1)$でも,定数$q$を同じ更新に含めよう.\footnote{$p=0$なら,$n\ge2$で$a_n=q$と直接分かる.$p=1$は後で扱う.行列による更新の表現自体は,どちらの場合も使える.}
状態を$\boldsymbol{x}_n=\begin{pmatrix}a_n\\1\end{pmatrix}$とすると,
\[
  \begin{aligned}
    A&=\begin{pmatrix}p&q\\0&1\end{pmatrix},\qquad
    \boldsymbol{x}_{n+1}=A\boldsymbol{x}_n,\\
    \boldsymbol{x}_n&=A^{n-1}\boldsymbol{x}_1.
  \end{aligned}
\]
元の\bookterm{recurrence}{漸化式}と,変わらない成分$1$をまとめたものである.

本編で$b_n=a_n-\alpha$と置いた操作は,この状態から
\[
  b_n=(1,-\alpha)\boldsymbol{x}_n
\]
という量を取り出すことに当たる.
実際,$\alpha=p\alpha+q$を使えば
\[
  \begin{aligned}
    (1,-\alpha)A
    &=(p,q-\alpha)\\
    &=(p,-p\alpha)=p(1,-\alpha).
  \end{aligned}
\]
本編の$b_{n+1}=pb_n$は,\bookterm{left-eigen}{左固有ベクトル}で取り出した量の更新だった.
状態の分解として見ると,$\alpha=p\alpha+q$から
\[
  A\begin{pmatrix}\alpha\\1\end{pmatrix}
  =
  \begin{pmatrix}\alpha\\1\end{pmatrix}.
\]
\bookterm{fixed}{不動点}は更新で変わらない方向になる.一方,
\[
  A\begin{pmatrix}1\\0\end{pmatrix}
  =
  p\begin{pmatrix}1\\0\end{pmatrix}.
\]
\bookterm{fixed}{不動点}からのずれは,一歩ごとに $p$ 倍される方向にある.
実際,初期状態は
\[
  \begin{pmatrix}a_1\\1\end{pmatrix}
  =
  \begin{pmatrix}\alpha\\1\end{pmatrix}
  +(a_1-\alpha)\begin{pmatrix}1\\0\end{pmatrix}
\]
と分けられる.$n-1$回更新して上の成分を読めば,
\[
  a_n=\alpha+(a_1-\alpha)p^{n-1}.
\]
本編の「\bookterm{fixed}{不動点}」と「ずれ」が,倍率の異なる二方向になった.

この\bookterm{eigenvalue}{固有値}$1,p$は
\BookMathLayout{\[
    \det(\lambda I-A)
    =
    \det\begin{pmatrix}\lambda-p&-q\\0&\lambda-1\end{pmatrix}
    =(\lambda-p)(\lambda-1)=0
  \]}{
  \[
    \begin{aligned}
      \det(\lambda I-A)&=\det\begin{pmatrix}\lambda-p&-q\\0&\lambda-1\end{pmatrix}\\
      &=(\lambda-p)(\lambda-1)=0
    \end{aligned}
  \]
}
からも得られる.本編の$\alpha=p\alpha+q$は不動点の\textbf{値},この\bookterm{characteristic}{特性方程式}は更新の\textbf{倍率}を求める式である.

本編では $p=1$ の場合を等差型として分けて扱った.
\bookterm{matrix}{行列}で書けば
\[
  A=\begin{pmatrix}1&q\\0&1\end{pmatrix},
  \qquad
  A^{n-1}=\begin{pmatrix}1&(n-1)q\\0&1\end{pmatrix}
\]
なので
\[
  a_n=a_1+(n-1)q
\]
となる.
$q\ne0$なら\bookterm{fixed}{不動点}はない.重なった\bookterm{eigenvalue}{固有値}$1$のもとで足し算が積み重なり,$n$が現れた.

\subsubsection{二回分の更新}\label{節：二回分の更新}
\sectionref*{節：一周期分をまとめる}では,$2x+1$の後に$3x+2$を計算し,$6x+5$を得た.
行列で書くと
\[
  A=\begin{pmatrix}2&1\\0&1\end{pmatrix},\qquad
  B=\begin{pmatrix}3&2\\0&1\end{pmatrix}.
\]
先に$A$,次に$B$なので,二回分は
\[
  BA=\begin{pmatrix}6&5\\0&1\end{pmatrix},\qquad
  BA\binom{x}{1}=\binom{6x+5}{1}
\]
となる.この例では$AB=BA$だが,一般には異なる.接続演習問4なら
\[
  U=\begin{pmatrix}-1&0\\0&1\end{pmatrix},\qquad
  V=\begin{pmatrix}2&1\\0&1\end{pmatrix}
\]
に対して,
\[
  VU=\begin{pmatrix}-2&1\\0&1\end{pmatrix},\qquad
  UV=\begin{pmatrix}-2&-1\\0&1\end{pmatrix}
\]
となる.更新を行う順序が,付加項の符号に現れた.
一般の$a_{n+1}=p_na_n+q_n$でも,$\binom{a_{n+1}}1=
  \begin{pmatrix}p_n&q_n\\0&1\end{pmatrix}\binom{a_n}1$と表せる.
一周期分は,先に行う\bookterm{matrix}{行列}を右に置いて掛け合わせる.途中の値は一回ずつの更新で求める.
\bookproblemref{practice-connections}{4}{接続演習問4}で本文との対応を確かめよう.


\subsubsection{分数を二つの成分で追う}\label{節：分数を二つの成分で追う}

\sectionref{節：一次分数型}で作った\bookterm{fixed}{不動点}からの差の比を,\bookterm{matrix}{行列}で見直す.
まず本編の
\[
  a_{n+1}=\frac{2a_n}{1-a_n}
\]
を使う.各項で$a_n\ne1$とし,$a_1\ne-1$とする.
一つの値を分子と分母に分けて持ち,$a_n=\dfrac{x_n}{y_n}$と書いてみよう.
$(x_1,y_1)=(a_1,1)$から
\[
  \begin{pmatrix}x_{n+1}\\y_{n+1}\end{pmatrix}
  =M\begin{pmatrix}x_n\\y_n\end{pmatrix},
  \qquad
  M=\begin{pmatrix}2&0\\-1&1\end{pmatrix}
\]
で更新すると,
\[
  \frac{x_{n+1}}{y_{n+1}}
  =\frac{2x_n}{y_n-x_n}
  =\frac{2a_n}{1-a_n}
\]
となり,元の更新に戻る.
本編で用いた$\displaystyle b_n=\frac{a_n}{a_n+1}$は,この二成分では
\[
  b_n=\frac{x_n}{x_n+y_n}
\]
である.分子は一回ごとに$2$倍されるが,分母は
\[
  x_{n+1}+y_{n+1}=2x_n+(y_n-x_n)=x_n+y_n
\]
と変わらない.従って,$b_{n+1}=2b_n$となる.
$x_1+y_1=a_1+1\ne0$なので,この比の分母はずっと零でない.

\bookterm{fixed}{不動点}$0,-1$を二成分で表すと,
\[
  M\begin{pmatrix}0\\1\end{pmatrix}
  =\begin{pmatrix}0\\1\end{pmatrix},
  \qquad
  M\begin{pmatrix}-1\\1\end{pmatrix}
  =2\begin{pmatrix}-1\\1\end{pmatrix}.
\]
対応する\bookterm{eigenvalue}{固有値}はそれぞれ$1,2$である.
差の比は,二方向の倍率の比を追っていた.

\bookstep{一般の一次分数型へ}
一般の係数でも,二成分の\bookterm{column-vector}{列ベクトル}で表せる.
\[
  a_{n+1}=\frac{pa_n+q}{ra_n+s}
\]
に対し,$r\ne0$, $ps-qr\ne0$を仮定し,各項で$ra_n+s\ne0$とする.
まず$(x_1,y_1)=(a_1,1)$とし,二つの数を
\[
  \begin{pmatrix}x_{n+1}\\y_{n+1}\end{pmatrix}
  =
  M\begin{pmatrix}x_n\\y_n\end{pmatrix},
  \qquad
  M=\begin{pmatrix}p&q\\r&s\end{pmatrix}
\]
で更新する.
$\displaystyle a_n=\frac{x_n}{y_n}$とすると
\[
  \frac{x_{n+1}}{y_{n+1}}
  =\frac{px_n+qy_n}{rx_n+sy_n}
  =\frac{pa_n+q}{ra_n+s}.
\]
また,$y_{n+1}=(ra_n+s)y_n$なので,もとの\bookterm{recurrence}{漸化式}の分母が零にならない限り,$y_n\ne0$が保たれる.
同じ零でない数で二成分を倍しても比は変わらない.この成分比を通して\bookterm{matrix}{行列}の作用を見る考え方を,射影変換という.

\bookterm{fixed}{不動点} $\alpha$ では
\[
  M\begin{pmatrix}\alpha\\1\end{pmatrix}
  =(r\alpha+s)\begin{pmatrix}\alpha\\1\end{pmatrix}.
\]
つまり\bookterm{fixed}{不動点}は\bookterm{eigenvalue}{固有値}ではなく,\bookterm{eigenvector}{固有ベクトル}の成分比である.
異なる二つの\bookterm{fixed}{不動点} $\alpha,\beta$ があり,対応する倍率を
$\lambda_\alpha=r\alpha+s,\ \lambda_\beta=r\beta+s$ とする.
状態を
\[
  \begin{pmatrix}x_n\\y_n\end{pmatrix}
  =c_n\begin{pmatrix}\alpha\\1\end{pmatrix}
  +d_n\begin{pmatrix}\beta\\1\end{pmatrix}
\]
と分けると,$\,c_n,d_n$ はそれぞれ $\lambda_\alpha,\lambda_\beta$ 倍される.
以下の比を取る際は$a_1\ne\beta$とする.\bookterm{fixed}{不動点}$\beta$を初項とする場合は$a_n=\beta$という定\bookterm{sequence}{数列}になる.
\bookterm{matrix}{行列}$M$は可逆で,$\lambda_\alpha\ne0$だから,$a_1\ne\beta$なら$c_n\ne0$,すなわち$a_n\ne\beta$が保たれる.
成分を比較すると
\[
  \begin{aligned}
    x_n&=\alpha c_n+\beta d_n,\qquad y_n=c_n+d_n,\\
    x_n-\alpha y_n&=(\beta-\alpha)d_n,\\
    x_n-\beta y_n&=(\alpha-\beta)c_n.
  \end{aligned}
\]
したがって,$\displaystyle a_n=\frac{x_n}{y_n}$より
\[
  \frac{a_n-\alpha}{a_n-\beta}=-\frac{d_n}{c_n},
  \qquad
  \frac{a_{n+1}-\alpha}{a_{n+1}-\beta}
  =\frac{\lambda_\beta}{\lambda_\alpha}
  \frac{a_n-\alpha}{a_n-\beta}.
\]
本編の\bookterm{fixed}{不動点}を求める方程式は$rx^2+(s-p)x-q=0$である.
解と係数の関係より$r(\alpha+\beta)=p-s$なので
\[
  \begin{aligned}
    \lambda_\beta&=r\beta+s=p-r\alpha,\\
    \lambda_\alpha&=r\alpha+s=p-r\beta.
  \end{aligned}
\]
よって公比は本編と同じ$\displaystyle \frac{\lambda_\beta}{\lambda_\alpha}=\frac{p-r\alpha}{p-r\beta}$になる.
置き換えは二つの固有方向の成分比を取り出す操作であり,最初の例は倍率$1,2$の場合だった.

\subsection{二つの値を更新する}\label{節：二つの値を更新する}

次は,二つの数を同時に更新する規則を考える.例えば
\[
  x_{n+1}=x_n+y_n,
  \qquad
  y_{n+1}=x_n
\]
という二つの規則は
\[
  \begin{pmatrix}x_{n+1}\\y_{n+1}\end{pmatrix}
  =
  \begin{pmatrix}1&1\\1&0\end{pmatrix}
  \begin{pmatrix}x_n\\y_n\end{pmatrix}
\]
と一つの\bookterm{matrix}{行列}で表せる.
初めを $(x_0,y_0)=(1,0)$ とすると,
\[
  (1,0),\ (1,1),\ (2,1),\ (3,2),\ldots
\]
と進み,$\,x_n=F_{n+1},\,y_n=F_n$ となる.
ここで $F_0=0,\,F_1=1$ はフィボナッチ\bookterm{sequence}{数列}である.
一歩前の$x$を$y$に残すので,$x$だけを見れば
\[
  x_{n+2}=x_{n+1}+x_n
\]
という二階の\bookterm{recurrence}{漸化式}になる.

\sectionref{節：隣接3項型}も,直前の二項を状態にすればよい.一般に
\[
  a_{n+2}=s a_{n+1}+t a_n
\]
に対して
\[
  \boldsymbol{v}_n=
  \begin{pmatrix}a_{n+1}\\a_n\end{pmatrix},
  \qquad
  B=\begin{pmatrix}s&t\\1&0\end{pmatrix}
\]
とすると,
\[
  B\boldsymbol{v}_n
  =
  \begin{pmatrix}s a_{n+1}+t a_n\\a_{n+1}\end{pmatrix}
  =
  \begin{pmatrix}a_{n+2}\\a_{n+1}\end{pmatrix}
  =\boldsymbol{v}_{n+1}.
\]
二行目が直前の項を残すので,$\boldsymbol{v}_{n+1}=B\boldsymbol{v}_n$という一階更新になる.
なお本編の形 $a_{n+2}+pa_{n+1}+qa_n=0$ で書く場合は,ここで $s=-p,\,t=-q$ とすればよい.

$B\begin{pmatrix}x\\y\end{pmatrix}=\lambda\begin{pmatrix}x\\y\end{pmatrix}$の二行目から$x=\lambda y$となる.
零でない\bookterm{eigenvector}{固有ベクトル}では$y\ne0$なので,$y=1$にそろえると$\begin{pmatrix}\lambda\\1\end{pmatrix}$になる.一行目から
\[
  s\lambda+t=\lambda^2,
  \qquad
  \lambda^2-s\lambda-t=0
\]
を得る.
これは$\lambda\ne0$なら,$a_n=\lambda^n$をもとの\bookterm{recurrence}{漸化式}に代入し,共通の冪を除いて得る\bookterm{characteristic}{特性方程式}と一致する.
$\lambda=0$も\bookterm{matrix}{行列}ではそのまま扱える.\sectionref*{節：行列の基本}で見たように,零の固有方向は初期状態にだけ寄与し,一回の更新で消える.
本編の符号で書けば $\lambda^2+p\lambda+q=0$ である.

この条件は\bookterm{linear}{線形}代数では
\BookMathLayout{\[
    \det(\lambda I-B)
    =
    \det\begin{pmatrix}\lambda-s&-t\\-1&\lambda\end{pmatrix}
    =\lambda^2-s\lambda-t=0
  \]}{
  \[
    \begin{aligned}
      \det(\lambda I-B)&=\det\begin{pmatrix}\lambda-s&-t\\-1&\lambda\end{pmatrix}\\
      &=\lambda^2-s\lambda-t=0
    \end{aligned}
  \]
}
とも書く.
本編の\bookterm{root}{特性根}は,更新行列の\bookterm{eigenvalue}{固有値}だった.

二つの異なる\bookterm{root}{特性根} $\lambda_1,\lambda_2$ があるとき,対応する\bookterm{eigenvector}{固有ベクトル}
\[
  \begin{pmatrix}\lambda_1\\1\end{pmatrix},
  \qquad
  \begin{pmatrix}\lambda_2\\1\end{pmatrix}
\]
は同じ直線上にはない.
初めの状態を二つの方向に分けておけば,$k$回の更新でそれぞれの部分が $\lambda_1^k,\,\lambda_2^k$ 倍される.
こうして\bookterm{general}{一般項}に\bookterm{root}{特性根}の冪が現れる.

さらに,本編で作った$a_{n+1}-\lambda_1a_n$という量も,同じ\bookterm{matrix}{行列}の中に見つかる.
解と係数の関係$\lambda_1+\lambda_2=s$, $\lambda_1\lambda_2=-t$を使うと
\[
  \begin{aligned}
    (1,-\lambda_1)B
    &=(s-\lambda_1,t)\\
    &=(\lambda_2,-\lambda_1\lambda_2)
    =\lambda_2(1,-\lambda_1).
  \end{aligned}
\]
したがって
\[
  \begin{aligned}
    b_n&=(1,-\lambda_1)\boldsymbol{v}_n
    =a_{n+1}-\lambda_1a_n,\\
    b_{n+1}&=\lambda_2 b_n.
  \end{aligned}
\]
根を入れ替えればもう一方も得られる.本編の$a_{n+2}=5a_{n+1}-6a_n$で作った$a_{n+1}-2a_n$は,$(1,-2)$で取り出す$3$倍の量だった.
\bookterm{eigenvector}{固有ベクトル}で状態を分ける方法と,\bookterm{left-eigen}{左固有ベクトル}で量を取り出す方法は,同じ更新の二通りの見方である.
二つの根が複素数でも,成分を複素数まで広げて同じ計算を行える.
実数の初期値から出発した場合,複素\bookterm{conjugate}{共役}な二方向の寄与が組になり,\bookterm{sequence}{数列}の項は実数のままになる.

\sectionref{節：連立型}の$A=\begin{pmatrix}2&1\\1&2\end{pmatrix}$でも,和と差$u_n=x_n+y_n$, $v_n=x_n-y_n$は
\[
  u_{n+1}=3u_n,\qquad v_{n+1}=v_n
\]
を満たす.前節の左固有ベクトル$(1,1)$は和を,$(1,-1)$は差を取り出していた.

図の$(1,1)$と$(1,-1)$の方向を使うと,座標軸を$45^\circ$回して見ることになる.
回すのは\textbf{座標の見方}であり,更新行列$A$自体が回転行列なのではない.
\BookFigSystem

向きをそろえて回転\bookterm{matrix}{行列}
\[
  Q=\frac{1}{\sqrt2}
  \begin{pmatrix}
    1&-1\\
    1&1
  \end{pmatrix}
\]
を用いると,
\[
  Q^{\mathsf T}
  \begin{pmatrix}x_n\\y_n\end{pmatrix}
  =\frac{1}{\sqrt2}
  \begin{pmatrix}u_n\\-v_n\end{pmatrix}
\]
となる.
$\boldsymbol z_n=\begin{pmatrix}x_n\\y_n\end{pmatrix}$とし,新しい座標を$\boldsymbol w_n=Q^{\mathsf T}\boldsymbol z_n$と書く.
直接掛けると$Q^{\mathsf T}Q=QQ^{\mathsf T}=I$だから,元へ戻す式は$\boldsymbol z_n=Q\boldsymbol w_n$である.
元の更新$\boldsymbol z_{n+1}=A\boldsymbol z_n$にこれを入れると,
\[
  \boldsymbol w_{n+1}
  =Q^{\mathsf T}\boldsymbol z_{n+1}
  =Q^{\mathsf T}AQ\boldsymbol w_n.
\]
元の座標へ戻す$Q$,更新する$A$,新座標へ直す$Q^{\mathsf T}$の順に作用している.
ここで
\[
  Q^{\mathsf T}AQ=
  \begin{pmatrix}3&0\\0&1\end{pmatrix}
\]
と表される.
だから,新座標では二つの等比型に分かれる.

実際に計算へ戻せば,
\[
  \begin{aligned}
    u_n&=3^{n-1}(x_1+y_1),\\
    v_n&=x_1-y_1,\\
    x_n&=\frac{u_n+v_n}{2},\qquad
    y_n=\frac{u_n-v_n}{2}
  \end{aligned}
\]
となる.
$x_1=1$, $y_1=0$では,前節の$A^k\begin{pmatrix}1\\0\end{pmatrix}$と$k=n-1$で一致する.
斜めの二方向を更新し,元の座標へ戻した結果である.

一般に連立\bookterm{recurrence}{漸化式}を
\BookMathLayout{\[
    \boldsymbol{z}_{n+1}=A\boldsymbol{z}_n,
    \qquad
    \boldsymbol{z}_n=\begin{pmatrix}a_n\\b_n\end{pmatrix},
    \qquad
    A=\begin{pmatrix}p&q\\r&s\end{pmatrix}
  \]}{
  \[
    \begin{aligned}
      \boldsymbol z_{n+1}&=A\boldsymbol z_n,\\
      \boldsymbol z_n&=\begin{pmatrix}a_n\\b_n\end{pmatrix},\quad
      A=\begin{pmatrix}p&q\\r&s\end{pmatrix}
    \end{aligned}
  \]
}
とし,零でない\bookterm{column-vector}{列ベクトル}$\boldsymbol{\ell}$を用いて$\,c_n=\boldsymbol{\ell}^{\mathsf T}\boldsymbol{z}_n$ とおく.
もし
\[
  \boldsymbol{\ell}^{\mathsf T}A
  =\gamma\boldsymbol{\ell}^{\mathsf T}
\]
なら,$\,c_{n+1}=\gamma c_n$ となる.
本編の連立型で $c_n=\alpha a_n+\beta b_n$ が\bookterm{geometric}{等比数列}となるよう $\alpha,\beta$ を選ぶ計算は,\bookterm{left-eigen}{左固有ベクトル}を求めていたのである.
この関係を満たす零でないベクトル$\boldsymbol{\ell}$が存在するための条件
\[
  \det(\gamma I-A)
  =\gamma^2-(p+s)\gamma+(ps-qr)=0
\]
は,本編で現れた\bookterm{characteristic}{特性方程式}と一致する.
異なる二つの\bookterm{eigenvalue}{固有値}があれば,対応する\bookterm{left-eigen}{左固有ベクトル}は比例しないので,取り出した二つの量から$a_n,b_n$を求め直せる.
二階の\bookterm{recurrence}{漸化式}も二成分の状態に直せば,この枠組みに入る.
フィボナッチ\bookterm{sequence}{数列}では $s=t=1$ なので
\[
  B=\begin{pmatrix}1&1\\1&0\end{pmatrix}
\]
であり,\bookterm{root}{特性根}は
\[
  \varphi=\frac{1+\sqrt5}{2},
  \qquad
  \psi=\frac{1-\sqrt5}{2}
\]
である.
初期値 $F_0=0,\,F_1=1$ に合わせて二つの方向の割合を決めると
\[
  F_n=\frac{\varphi^n-\psi^n}{\sqrt5}
\]
を得る.
$|\psi|<1<\varphi$だから,$\psi^n$は符号を交互に変えながら減衰し,$\varphi^n$が増大を支配する.
\bookterm{root}{特性根}は長期的な動きも表す.

\bookterm{characteristic}{特性方程式}が重解 $r\ne0$ を持つ場合には,二つの異なる倍率の方向に分けられない.
\bookterm{general}{一般項}に $n$ が現れる理由は,\bookterm{recurrence}{漸化式}に直接戻ると見やすい.
重解のときは
\[
  a_{n+2}=2r a_{n+1}-r^2a_n
\]
だから,$\,a_n=r^n b_n$ と置いて代入すると
\[
  b_{n+2}-2b_{n+1}+b_n=0
\]
となる.
差が一定だから$b_n=C+Dn$,つまり$a_n=(C+Dn)r^n$となる.
\sectionref{節：付加項がない場合}の重解の例題で得た$a_n=n2^{n-1}$も,等比的な倍率$2^{n-1}$を取り除くと$n$という等差的な増加が残る例である.
\sectionref{節：階差等比複合型}の指数型で倍率が一致した場合とも,等比部分を取り除いて差を取る工夫が共通する.
たとえば重解が $1$ のときは
\BookMathLayout{\[
    a_{n+2}-2a_{n+1}+a_n
    =(a_{n+2}-a_{n+1})-(a_{n+1}-a_n)=0
  \]}{
  \[
    \begin{aligned}
      a_{n+2}-2a_{n+1}+a_n\\
      {}=(a_{n+2}-a_{n+1})-(a_{n+1}-a_n)\\
      {}=0
    \end{aligned}
  \]
}
より,隣り合う項の差が一定で $a_n=C+Dn$ となる.
重解が $r=0$ の場合は $a_{n+2}=0$ となる退化した\bookterm{recurrence}{漸化式}であり,直接調べればよい.

\subsection{角度の整数倍と多項式}\label{節：角度の整数倍と多項式}

\sectionref{節：特性根が虚数の場合}では,複素数の冪と正弦・余弦という\bookterm{basis}{基底}を使い分けた.
\sectionref{節：三角関数型}では,値の更新を角度の更新に直した.
この二つを「整数倍の角を表す多項式」でつなぐ.

まず,一定の角$\theta$だけ回す操作を考える.
\bookterm{addition}{加法定理}より
\[
  R=\begin{pmatrix}\cos\theta&-\sin\theta\\
    \sin\theta&\cos\theta\end{pmatrix}
\]
とおけば
\[
  \begin{pmatrix}\cos((n+1)\theta)\\\sin((n+1)\theta)\end{pmatrix}
  =R\begin{pmatrix}\cos(n\theta)\\\sin(n\theta)\end{pmatrix}.
\]
この\bookterm{matrix}{行列}の\bookterm{characteristic}{特性方程式}は
\[
  \lambda^2-2\cos\theta\,\lambda+1=0
\]
であり,$0<\theta<\pi$なら\bookterm{eigenvalue}{固有値}は$\cos\theta\pm i\sin\theta$である.
複素数の冪による回転が,\bookterm{eigenvalue}{固有値}にも現れた.

横座標だけを追うと,余弦の\bookterm{addition}{加法定理}より
\[
  a_{n+2}=2\cos\theta\,a_{n+1}-a_n
\]
という隣接3項間\bookterm{recurrence}{漸化式}になる.
ここでは\textbf{一定角を足す}.本編の\bookterm{trigonometric}{三角関数}型の\textbf{角度を二倍にする}更新とは異なるが,以下の多項式はどちらの整数倍角も表せる.

$x=\cos\theta$とすると,\bookterm{addition}{加法定理}から
\[
  \cos((n+1)\theta)=2x\cos(n\theta)-\cos((n-1)\theta).
\]
同じ規則で多項式の列を
\begin{align*}
  &T_0(x)=1,\qquad T_1(x)=x,\\
  &T_{n+1}(x)=2xT_n(x)-T_{n-1}(x)
  \qquad(n\ge1)
\end{align*}

と定める.
最初から足し算と掛け算だけで作るので,すべて多項式になる.これを\bookterm{chebyshev}{チェビシェフ多項式}という.

この定義から
\[
  T_2(x)=2x^2-1,\qquad
  T_3(x)=4x^3-3x
\]
が得られる.
二倍角・三倍角の公式と同じ形である.一般にも,$n=0,1$の値と更新規則が一致するので,帰納法から
\[
  T_n(\cos\theta)=\cos(n\theta).
\]
整数倍の角の余弦を,$x=\cos\theta$の多項式で表せた.図は$T_2$による二倍角の横座標を示す.
\BookFigChebyshev

ここで本編の$\displaystyle a_1=\frac{\sqrt3}{2}$, $a_{n+1}=2a_n^2-1$に戻ろう.
本編では$\displaystyle a_1=\cos(\frac{\pi}{6})$から角度を二倍ずつ進めたが,多項式の言葉では
\[
  a_{n+1}=T_2(a_n),\qquad
  a_n=T_{2^{n-1}}(a_1)
\]
と書ける.
初項そのものの多項式で\bookterm{general}{一般項}を表せた.その理由は次の合成の関係にある.

一般に,$x=\cos\theta$なら
\BookMathLayout{\[
    T_m(T_k(x))
    =T_m(\cos(k\theta))
    =\cos(mk\theta)
    =T_{mk}(x).
  \]}{
  \[
    \begin{aligned}
      T_m(T_k(x))&=T_m(\cos(k\theta))\\
      &=\cos(mk\theta)=T_{mk}(x).
    \end{aligned}
  \]
}
この等式は $[-1,1]$ のすべての $x$ で成り立ち,両辺が多項式なので,多項式の恒等式として実数全体で成り立つ.
したがって$T_2$を繰り返し合成する操作は,添字を二倍ずつ進める$T_2,T_4,T_8,\ldots$にまとまる.
本編の「整数倍の公式へ」で扱った$m$倍角にも同じことがいえる.
整数$m\ge2$に対して$a_{n+1}=T_m(a_n)$なら
\[
  a_n=T_{m^{n-1}}(a_1)
\]
である.
$-1\le a_1\le1$なら$a_1=\cos\theta$と表せるため,これは本編の$a_n=\cos(m^{n-1}\theta)$と一致する.
多項式としての式は,\bookterm{trigonometric}{三角関数}で表せない初項についても成り立つ.

多項式列を\bookterm{recurrence}{漸化式}で作る例は,\bookterm{chebyshev}{チェビシェフ多項式}だけではない.
\bookterm{legendre}{ルジャンドル多項式} $P_n(x)$ は,初期値
\[
  P_0(x)=1,\qquad P_1(x)=x
\]
と三項\bookterm{recurrence}{漸化式}
\BookMathLayout{
  \[
    \begin{aligned}
      (n+1)P_{n+1}(x)&=(2n+1)xP_n(x)\\
      &\quad-nP_{n-1}(x)\qquad(n\ge1)
    \end{aligned}
  \]
}{
  \[
    \begin{aligned}
      (n+1)P_{n+1}(x)&=(2n+1)xP_n(x)\\
      &\quad-nP_{n-1}(x)\qquad(n\ge1)
    \end{aligned}
  \]
}
で定まる.
例えば
\[
  P_2(x)=\frac{3x^2-1}{2}
\]
である.
\bookterm{chebyshev}{チェビシェフ多項式}と同じく三項\bookterm{recurrence}{漸化式}から作られるが,\bookterm{legendre}{ルジャンドル多項式}では係数が $n$ にも依存する.

\bookterm{legendre}{ルジャンドル多項式}は,区間 $[-1,1]$ 上で重み $1$ に関して直交する.
ここでいう直交性とは,異なる次数の二つを掛けて区間全体で積分すると $0$ になる性質である.
\[
  \int_{-1}^{1}P_m(x)P_n(x)\,dx=0
  \qquad(m\ne n)
\]
この性質は関数の表示や近似に役立つ.ここでは,三項\bookterm{recurrence}{漸化式}が生む別の多項式列として紹介する.

\subsection*{コラム：有限個の点から多項式を作る}

最初の三項が $1,2,3$ なら,
$x$ と $x+(x-1)(x-2)(x-3)$ はどちらもその三項に一致するが,
$n=4$ ではそれぞれ $4$ と $10$ になる.
有限個の値が一致しても,その後の値まで同じとは限らない.
与えられた値を通る多項式を,一般の場合にも作ってみよう.

最初の$N$項$a_1,\ldots,a_N$に対して,
\[
  L(x)=\sum_{k=1}^{N}a_k
  \prod_{\substack{1\le j\le N\\j\ne k}}
  \frac{x-j}{k-j}
\]
と置くと, $L(k)=a_k$ $(1\le k\le N)$ となる.
$x=k$ のとき $k$ 番目の積は $1$ になり,ほかの積には $x-k$ の因子があるためである.
次数を $N-1$ 以下に制限した多項式の中では,この $L$ は一意である.二つの候補の差は $1,\ldots,N$ を根にもつ次数 $N-1$ 以下の多項式なので,零多項式になる.

次数制限はデータからは決まらない.$L(x)+c\prod_{j=1}^{N}(x-j)$も,任意の定数$c$で最初の$N$個の値を保つ.
次の図は,別の三点を通り,途中の形が異なる二つの曲線を描いた例である.
\BookFigLagrange

\bookterm{interpolation}{ラグランジュ補間}は一つの延長候補を作る.それを元の規則と結論づけるには,次数や\bookterm{recurrence}{漸化式}など追加の条件が要る.

\subsection{不動点付近の線形化}\label{節：不動点付近の線形化}

\sectionref*{節：特性方程式型}では,不動点からのずれが$p$倍され,$|p|<1$なら近づいた.
\bookterm{nonlinear}{非線形}の更新でも,微分係数と一次近似で不動点付近の倍率を読める.
一項前から次が決まる\bookterm{recurrence}{漸化式}
\[
  a_{n+1}=f(a_n)
\]
を考え,$\,\alpha$ を\bookterm{fixed}{不動点} $f(\alpha)=\alpha$ とする.
そこからのずれを $e_n=a_n-\alpha$ と置くと
\[
  e_{n+1}=f(\alpha+e_n)-f(\alpha)
\]
である.
$f$ が $\alpha$ で微分可能なら,微分係数の定義から
\[
  \frac{f(\alpha+e)-f(\alpha)}{e}\longrightarrow f'(\alpha)
  \qquad(e\longrightarrow0)
\]
である.
つまり,ずれ $e$ が十分小さい範囲では
\[
  f(\alpha+e)-f(\alpha)\approx f'(\alpha)e
\]
となる.
不動点の近くで曲線を接線に置き換えたのである.収束も調べられる.
もし $|f'(\alpha)|<1$ なら,$\,|f'(\alpha)|$ より大きく $1$ より小さい数 $c$ を選べる.
微分係数の定義により,$\,a_n$ が $\alpha$ に十分近ければ
\[
  |e_{n+1}|
  =|f(\alpha+e_n)-f(\alpha)|
  \le c|e_n|
\]
となる.
誤差が毎回$c$倍以下に縮むので,十分近い初期値は近くにとどまり,$\alpha$へ収束する.
行列の\bookterm{eigenvalue}{固有値}と同じく,$f'(\alpha)$が局所的なずれの倍率を表す.
一次分数型で先ほどの二つの異なる\bookterm{fixed}{不動点}がある場合には
\[
  \begin{aligned}
    f'(\alpha)
    &=\frac{ps-qr}{(r\alpha+s)^2}\\
    &=\frac{\lambda_\alpha\lambda_\beta}{\lambda_\alpha^2}
    =\frac{\lambda_\beta}{\lambda_\alpha}.
  \end{aligned}
\]
特性方程式から$ps-qr=\lambda_\alpha\lambda_\beta$である.
比の置き換えで得た\textbf{正確な公比}と,微分による\textbf{局所的な倍率}が一致した.
\bookterm{linear}{線形}代数と微分は,一歩の変化を倍率で測る見方でつながる.

\subsection{整数から作る近似分数}\label{節：整数から作る近似分数}
\sectionref*{節：ニュートン法}では,$a_1=2$から$\sqrt2$を近似した.
ここでは同じ更新式$a_{n+1}=\frac12\left(a_n+\dfrac{2}{a_n}\right)$で,初項を$a_1=1$に変えてみる.
\[
  a_2=\frac32,\qquad a_3=\frac{17}{12},\qquad a_4=\frac{577}{408}
\]
となる.初項$1,2$のどちらも$\dfrac{3}{2}$へ移るため,第2項以降は本編と一致する.
分子と分母を追い,連立型・一次分数型・\bookterm{newton}{ニュートン法}をつなげる.微積分は使わない.

\bookstep{分子と分母の関係}
$3^2-2\cdot2^2=17^2-2\cdot12^2=577^2-2\cdot408^2=1$である.
次の\bookterm{sequence}{数列}を作ろう.
\[
  x_k+y_k\sqrt2=(3+2\sqrt2)^k\quad(k\ge1)
\]
$3+2\sqrt2$倍して係数を比べると,
\[
  x_{k+1}=3x_k+4y_k,\qquad y_{k+1}=2x_k+3y_k
\]
となる.$x_1=3$, $y_1=2$なので,両方とも正の整数である.
\bookterm{conjugate}{共役}な式との積から
\[
  x_k^2-2y_k^2=(3+2\sqrt2)^k(3-2\sqrt2)^k=1
\]
が保たれる.整数の式を保つ更新は,
\bookproblemref{practice-bridge}{9}{発展入門問9}の見方とも共通する.

\bookstep{比から一次分数型へ}
$r_k=\frac{x_k}{y_k}$とすると,連立の更新から
\[
  r_{k+1}=\frac{3r_k+4}{2r_k+3}
\]
である.分母は正であり,\bookterm{fixed}{不動点}を求めると$\pm\sqrt2$になる.
一次分数型で用いた差の比は,
\[
  \frac{r_{k+1}-\sqrt2}{r_{k+1}+\sqrt2}
  =(3-2\sqrt2)^2
  \frac{r_k-\sqrt2}{r_k+\sqrt2}
\]
となる.整数の組とその比は,同じ更新の二つの表現である.

\bookstep{例：$\dfrac{17}{12}$の次の近似分数}
本編で使った更新を$\dfrac{17}{12}$に施すと,
\[
  \frac12\left(\frac{17}{12}+\frac{24}{17}\right)=\frac{577}{408}
\]
となる.一方,整数の更新を一回進めた$r_3$は$\dfrac{99}{70}$である.
ニュートン法は,この列で次の分数へ一つ進んでいるわけではない.

\bookstep{別解：ニュートン法と添字の倍増}
$(x_k+y_k\sqrt2)^2$の展開から,
\[
  x_{2k}=x_k^2+2y_k^2,\qquad y_{2k}=2x_ky_k
\]
なので,
\[
  r_{2k}=\frac12\left(r_k+\frac2{r_k}\right).
\]
この\bookterm{newton}{ニュートン法}型の式から
\[
  a_n=r_{2^{n-2}}\qquad(n\ge2)
\]
となる.$\dfrac{3}{2}$, $\dfrac{17}{12}$, $\dfrac{577}{408}$は$r_1$, $r_2$, $r_4$であり,連続する三項ではない.
\BookFigPellDoubling
$x_k^2-2y_k^2=1$を因数分解すれば,
\[
  r_k-\sqrt2=\frac1{y_k^2(r_k+\sqrt2)}>0.
\]
$y_{k+1}>3y_k$なので$y_k\to\infty$であり,$r_k\to\sqrt2$となる.
添字の倍増からも,ニュートン法の速さを読める.

\bookstep{ペル方程式}
平方数でない正の整数$d$に対し,$x^2-dy^2=1$の整数解を求める問題を\textbf{\bookterm{pell}{ペル方程式}}という.\footnote{ここでは$d=2$の具体例を扱い,一般の解の存在や全解の分類には進まない.}

\bookstep{連分数にも現れる分数}
$\sqrt2=1+\frac1{1+\sqrt2}$から,有限段階の近似を
\[
  C_0=1,\qquad C_{j+1}=1+\frac1{1+C_j}
  \quad(j\ge0)
\]
で定める.$C_1=\dfrac{3}{2}$, $C_2=\dfrac{7}{5}$, $C_3=\dfrac{17}{12}$である.
$C_j=\dfrac{P_j}{Q_j}$とすると,$P_0=Q_0=1$から
\[
  P_{j+1}=P_j+2Q_j,\qquad Q_{j+1}=P_j+Q_j
\]
が得られ,$P_j+Q_j\sqrt2=(1+\sqrt2)^{j+1}$となる.
$(1+\sqrt2)^2=3+2\sqrt2$だから,
\[
  C_{2k-1}=r_k\qquad(k\ge1)
\]
である.この有限段階の値を\textbf{\bookterm{continued}{連分数}の近似分数}と呼ぶ.
偶数番号では$P_j^2-2Q_j^2=-1$,奇数番号では$1$となり,
同じ誤差の因数分解で$C_j\to\sqrt2$も確認できる.

\begin{bookpoint}
  \textbf{三つの解法のつながり}\par
  分子・分母を追えば連立型,比を追えば一次分数型,
  添字を倍に進めれば平方根の\bookterm{newton}{ニュートン法}になる.
  追う量と更新回数を変えると,解法がつながる.
\end{bookpoint}
保存される式と誤差を\bookproblemref{practice-connections}{6}{接続演習問6}で確かめよう.
\bookterm{pell}{ペル方程式}と\bookterm{continued}{連分数}の補足資料は,\chapterref*{章：参考文献}のKeith Conradの二つの解説にある.


  % 付録：階差型・総和型から微積分と近似計算へ

\section{和と差から微積分へ}\label{節：和と差から微積分へ}

\sectionref{節：階差型}は差から元の\bookterm{sequence}{数列}を求め,\sectionref{節：総和型}は和から差を取り出した.
この逆向きの二つの操作を,微分と積分につなげよう.
隣の項との差は,一歩進んだときの変化である.
変化を取り出す操作と,その変化を足して元へ戻す操作を,同じ流れの中で見ていく.

\subsection{和と差の対応}
以下では,最初の項を$a_0$と書くことがある.
例えば平方の列$0,1,4,9,\ldots$から隣り合う差を取ると,$1,3,5,\ldots$になる.
\bookterm{sequence}{数列}$a_0,a_1,a_2,\ldots$の隣り合う項の差
\[
  \Delta a_n=a_{n+1}-a_n
\]
を一\bookterm{difference}{階差}分という.
本編の階差に$\Delta$という記号と「\bookterm{difference-operator}{差分}」の名前を付けた.
たとえば$a_n=n^2$なら
\[
  \Delta a_n=2n+1,\qquad \Delta^2a_n=2.
\]
二次式,一次式,定数の順になる.$\Delta^2a_n=\Delta a_{n+1}-\Delta a_n$である.

\sectionref*{節：特性方程式型}の別解でも,添字を一つずらして引くことで定数項$q$を消し,
\[
  \Delta a_{n+1}=p\,\Delta a_n
\]
という等比型を作った.
これも\bookterm{difference-operator}{差分}で扱いやすい量を取り出す例である.

逆に,\bookterm{difference-operator}{差分}を足し合わせれば元の\bookterm{sequence}{数列}を求め直せる.
\[
  \begin{aligned}
    a_n-a_0
    &=(a_1-a_0)+\cdots+(a_n-a_{n-1})\\
    &=\sum_{k=0}^{n-1}\Delta a_k.
  \end{aligned}
\]
途中の項が消える.本編の\bookterm{difference}{階差}型の和は離散的な積分,総和型の$S_{n+1}-S_n=a_{n+1}$はその逆になる.図で微積分との対応を見る.
\BookFigDifference

この見方は,\sectionref{節：階差等比複合型}で引く多項式$g(n)$を探したことにもつながる.
特に$p=1$なら,$g(n+1)-g(n)=f(n)$を満たす式を探す問題になる.
たとえば$S_n=\sum_{k=1}^n k^2$については
\[
  S_n-S_{n-1}=n^2,\qquad S_0=0.
\]
二次式の差を逆向きに戻すので,三次式
$S_n=An^3+Bn^2+Cn+D$を候補にする.
その差は
\[
  S_n-S_{n-1}
  =3An^2+(-3A+2B)n+(A-B+C).
\]
係数比較から$\displaystyle A=\frac{1}{3}$, $\displaystyle B=\frac{1}{2}$, $\displaystyle C=\frac{1}{6}$となり,$S_0=0$から$D=0$である.
同じ初期値と差を満たすので
\[
  S_n=\frac{n(n+1)(2n+1)}6
\]
が確かめられる.
和の公式は,差を取ると目的の式になるものを探して作れる.

\subsection{二項係数で多項式を表す}\label{節：二項係数で多項式を表す}
平方和を,今度は\sectionref*{節：二項定理}の二項係数で求める.
ここでは$n,k,N$は非負整数とし,
\[
  \binom nk=\frac{n!}{k!(n-k)!}\quad(0\le k\le n)
\]
とする.$0!=1$, $k>n$では$\binom nk=0$と約束する.

\bookstep{例題：二項係数で平方和を求める}
$n^2=n(n-1)+n$だから,
\[
  n^2=2\binom n2+\binom n1
\]
と表せる.
部品の差は
\[
  \binom{n+1}{3}-\binom n3=\binom n2\quad(n\ge0)
\]
となる.$n\ge3$では階乗表示を展開して確かめられ,$n=0,1,2$も上の規約で成り立つ.
同じように,$\binom{n+1}{2}-\binom n2=\binom n1$である.
$j=0$から$N-1$まで足すと途中が消え,従って
\BookMathLayout{\[
    \sum_{j=0}^{N-1}j^2=2\binom N3+\binom N2
    =\frac{N(N-1)(2N-1)}6.
  \]}{\[
    \begin{aligned}
      \sum_{j=0}^{N-1}j^2&=2\binom N3+\binom N2\\
      &=\frac{N(N-1)(2N-1)}6.
    \end{aligned}
  \]}
$N=n+1$とすれば,先ほどの$1^2+\cdots+n^2$の公式に戻る.
二項係数なら,添字を一つ下げるだけで差を表せる.

\bookstep{一般の添字で確かめる}
$n+1$個から$k$個を選ぶとき,特定の一個を選ぶかどうかで分ければ,
\[
  \binom{n+1}{k}=\binom nk+\binom n{k-1}\quad(k\ge1)
\]
である.従って
\[
  \Delta\binom nk=\binom n{k-1}\quad(k\ge1)
\]
となる.$\Delta$は前進の差で,$\binom n0=1$の差は別に零となる.

逆に差を足すと,途中の項が消えるので,
\[
  \sum_{j=0}^{N-1}\binom jk=\binom N{k+1}
  \quad(N,k\ge0)
\]
となる.$N=0$では空和$0$,$k=0$では和が$N$となり,右辺と一致する.
\BookFigBinomialDifference

\begin{bookpoint}
  \textbf{式の部品を選ぶ}\par
  目的の操作が簡単になる\bookterm{basis}{基底}を選ぶ.
  \sectionref*{節：隣接3項型}の見方が,多項式の表示にも現れた.
\end{bookpoint}
実際,$\binom n0,\ldots,\binom nd$はそれぞれ次数が$0,\ldots,d$で,
最高次係数が非零である.高い次数から係数を合わせれば,
$d$次以下の多項式を一意にその\bookterm{combination}{線形結合}で表せる.
三乗和は\bookproblemref{practice-connections}{7}{接続演習問7}で試そう.


\subsection{二方向の階差}\label{節：二方向の階差}
非負整数$n,m$で数表の値$a_{n,m}$を表す.
\sectionref*{節：階差型}の和と差を,横と縦の二方向で使おう.

\bookstep{例：距離の二乗と階差}
原点から点$(n,m)$までの距離は
\[ a_{n,m}=\sqrt{n^2+m^2} \]
である.
根号を取り除くため$b_{n,m}=a_{n,m}^2$とすると,
一方向の差は$b_{n+1,m}-b_{n,m}=2n+1$となり,$m$によらない.
従って,さらにもう一方向へ差を取れば零になる.

\bookstep{零の差から一定の差へ}
この計算を広げ,境界の値
\[
  b_{n,0}=n^2,\qquad b_{0,m}=m^2
\]
と,一定の実数$c$について
\BookMathLayout{\[
    b_{n+1,m+1}-b_{n+1,m}-b_{n,m+1}+b_{n,m}=c
  \]}{\[
    \begin{aligned}
      b_{n+1,m+1}-b_{n+1,m}\\
      {}-b_{n,m+1}+b_{n,m}=c
    \end{aligned}
  \]}
とする.$b_{n+1,m+1}$は添字の和が小さい三つの値から決まるので,$n+m$についての帰納法で境界から一意に求まる.

横方向の差$d_{n,m}=b_{n+1,m}-b_{n,m}$を追うと,
\[
  d_{n,m+1}-d_{n,m}=c,\qquad d_{n,0}=2n+1
\]
なので$d_{n,m}=2n+1+cm$である.これを横に足し戻せば,
\[
  b_{n,m}=n^2+m^2+cnm
\]
となる.$c$が一定なので,\textbf{二方向の\bookterm{difference}{階差}を逆順に足し戻せた}.

\bookstep{数表と図形}
二つの単位ベクトル$\boldsymbol u,\boldsymbol v$のなす角を$\theta$とすると,
$c=2\cos\theta$の場合には
\[
  b_{n,m}=|n\boldsymbol u+m\boldsymbol v|^2
\]
と読める.数表なら$c$は任意の実数でよいが,この距離の二乗という解釈には$-2\le c\le2$が必要になる.

\bookstep{例：格子経路を数える}
一方,格子点$(0,0)$から$(n,m)$へ,右と上へ一歩ずつ進む道の数を$A_{n,m}$とする.
最後の一歩を分ければ,
\[
  A_{n,m}=A_{n-1,m}+A_{n,m-1}\quad(n,m\ge1)
\]
であり,境界では$A_{n,0}=A_{0,m}=1$である.
\BookFigLatticeRecurrence
下の数表は,上と左の値を足して作ったものである.
\[
  \begin{array}{c|rrrr}
    m\backslash n&0&1&2&3\\\hline
    0&1&1&1&1\\
    1&1&2&3&4\\
    2&1&3&6&10\\
    3&1&4&10&20
  \end{array}
\]
全$n+m$歩から上へ進む$m$歩の位置を選ぶので,
\[
  A_{n,m}=\binom{n+m}{m}
\]
である.最後の一歩の分類は,二項係数の加法公式でもある.
境界と帰納法による確認は\bookproblemref{practice-connections}{3}{接続演習問3}へ.

\bookstep{二つの例の共通点}
このように二つの添字をもつものを\textbf{\bookterm{double}{二重数列}}という.
一変数で使った階差と帰納法を二方向へ広げた見方である.

\begin{bookpoint}
  \textbf{二つの添字と初期条件}\par
  \bookterm{initial}{初期条件}となる境界と更新順序を確かめる.
  二方向の規則を別々に与えるなら,同じ位置を二通りに計算して一致する必要がある.
\end{bookpoint}
格子経路と二項係数の資料は参考文献のNIST DLMF §26.3へ.


\subsection{差分商と微分}
間隔$h>0$で平方を記録すると,$y(t)=t^2$について
\[
  \frac{y(t+h)-y(t)}h
  =\frac{(t+h)^2-t^2}h=2t+h.
\]
となり,$h\to0$で傾き$2t=y'(t)$に近づく.一般に$t_n=nh$での値を$a_n$とすると,単位時間あたりの変化は
\[
  \frac{a_{n+1}-a_n}{h}
\]
である.
微分可能な関数$y$について,\textbf{時刻$t$を固定}して$h$を小さくすると,\footnote{基点も動かす近似では,例えば$y'$がある閉区間で連続なら,その区間に収まる$t,t+h$について差分商と$y'(t)$の差が一様に小さくなる.ここで使う$y(t)=t^2$はこの条件を満たす.}
\[
  \frac{y(t+h)-y(t)}h\longrightarrow y'(t)
\]
となる.格子点$t_n$での傾きもこの比で近似できるが,基点も動かす場合の一様な近似には追加条件が要る.

逆に,微分方程式$y'=F(t,y)$を数値で解くには,今の点での傾きから
\[
  a_{n+1}=a_n+hF(t_n,a_n)
\]
と次の値を予測できる.
これをオイラー法という.接線の方向へ一歩ずつ進む\bookterm{recurrence}{漸化式}で曲線を近似する.

例えば$y'=-y$, $y(0)=1$の解は$e^{-t}$で,$0$に近づく.
オイラー法では
\[
  a_{n+1}=(1-h)a_n,\qquad a_0=1
\]
だから,$n\ge1$では$a_n=(1-h)^n$である.最初の値は$a_0=1$のまま保つ.
本編の等比型で調べたように,$0$への収束条件は$|1-h|<1$, すなわち$0<h<2$になる.
$0<h<1$では正のまま減少し,$h=1$なら一回で$0$になる.
$1<h<2$では符号を交互に変えながら近づくが,$h=2$では$1,-1,1,-1,\ldots$と繰り返して収束しない.
さらに$h>2$では絶対値が増大する.
本当の解が減衰しても,刻み幅次第で近似の列は増大する.等比型の公比は,数値計算の安定性も教える.

また,正の時刻$t$を$N$等分して$\displaystyle h=\frac{t}{N}$とすれば,時刻$t$における近似値は
\[
  a_N=\left(1-\frac{t}{N}\right)^N
  \longrightarrow e^{-t}.
\]
この例では,刻みを細かくした近似の収束を指数関数の\bookterm{limit}{極限}で確かめられる.一般の保証には関数や刻み幅の条件が要る.

\subsection{ニュートン法の誤差}
\sectionref{節：ニュートン法}では,\BookMathLayout{}{\\}%
$a_{n+1}=\frac12\left(a_n+\dfrac{2}{a_n}\right)$の$\sqrt2$への誤差がほぼ二乗された.
一般の\bookterm{newton}{ニュートン法}でも,誤差が二乗の定数倍以下になる条件を調べよう.

方程式$g(x)=0$に対する\bookterm{newton}{ニュートン法}の更新は
\[
  x_{n+1}=x_n-\frac{g(x_n)}{g'(x_n)}
\]
である.
$g(\alpha)=0$, $g'(\alpha)\ne0$とし,根$\alpha$の近くで$g$は二回連続微分可能とする.
本編のずれの置き換えと同じく,$e_n=x_n-\alpha$とおくと
\[
  e_{n+1}
  =\frac{g'(x_n)e_n-g(x_n)}{g'(x_n)}.
\]
分子は,接線による変化と実際の変化との差である.
二階微分$g''$の絶対値がこの近くで$M$以下なら
\[
  \begin{aligned}
    g'(x_n)e_n-g(x_n)
    &=\int_\alpha^{x_n}\{g'(x_n)-g'(t)\}\,dt,\\
    |g'(x_n)e_n-g(x_n)|
    &\le \frac{M}{2}|e_n|^2.
  \end{aligned}
\]
平均値の定理で$|g'(x_n)-g'(t)|\le M|x_n-t|$と評価して積分した.$x_n<\alpha$でも絶対値で同じ評価になる.
また,$g'(\alpha)\ne0$と連続性により,十分近くでは$|g'(x_n)|\ge m$となる正の定数$m$を選べる.
したがって
\[
  |e_{n+1}|\le\frac{M}{2m}|e_n|^2.
\]
十分小さな誤差は二乗の定数倍以下に縮むので,近くにとどまり急速に根へ収束する.
本編の$\sqrt2$の例で得た正確な式
\[
  e_{n+1}=\frac{e_n^2}{2a_n}
\]
は,その具体例である.

前節の\bookterm{fixed}{不動点}の見方でも確かめられる.
更新関数$\displaystyle N(x)=x-\frac{g(x)}{g'(x)}$について
\[
  N'(\alpha)
  =\left.\frac{g(x)g''(x)}{g'(x)^2}\right|_{x=\alpha}=0.
\]
つまり,根のところでは一次のずれが消え,次の誤差は二乗の定数倍以下になる.
二次の項まで消えると,さらに速く近づく場合がある.\footnote{例えば$g(x)=x+x^3$の根$0$では,$N(x)=\dfrac{2x^3}{1+3x^2}$となり,誤差は三次の大きさになる.$g(x)=x$なら一回で根に移る.}
等比型の一定倍率と違い,\bookterm{newton}{ニュートン法}では誤差が小さいほど残る割合も小さくなる.

接線の傾きを毎回微分で求めるのが難しいなら,過去二点を結ぶ割線の傾きで代用できる.
その更新は
\[
  x_{n+1}
  =x_n-g(x_n)\frac{x_n-x_{n-1}}{g(x_n)-g(x_{n-1})}
\]
となり,割線法と呼ばれる.
分母が零でない範囲で用いる.
直前の二項を状態にする点は前節と共通する.\bookterm{recurrence}{漸化式}は,近似の規則を作り収束を調べる道具にもなる.

  % BOOK-FILE-ID: f-71d7afa99e82
% 付録：本編の式変形を母関数と数列の性質へ広げる

\section{数列をまとめて見る}\label{節：数列をまとめて見る}

本編では,添字をずらして引く操作を何度も使ってきた.
ここではそれを,少し突飛な方法で見直してみよう.
数列の各項を,$x$の累乗の係数として1つの式に並べてしまうのだ.
添字は項の位置を表していたが,その位置を$x$の次数として記録すれば,$x$を掛けるだけで1つ先へずらせる.
等比型$a_0=1$, $a_{n+1}=2a_n$の列$1,2,4,8,\ldots$なら,
\[
  G(x)=1+2x+4x^2+8x^3+\cdots
\]
と書く.ここでは$x$に値を入れて和を求めるのではなく,各次数の係数を記録している.
$2x$を掛けると,一つ先へずれた係数が元の係数と重なる.
\[
  \begin{aligned}
    2xG(x)&=2x+4x^2+8x^3+\cdots,\\
    G(x)-2xG(x)&=1.
  \end{aligned}
\]
したがって$G(x)=\frac{1}{1-2x}$である.「ずらして引く」が,ただの掛け算と引き算になったのだ.
このように各項を係数として並べた式を\bookterm{generating}{母関数}といい,本編の\bookterm{characteristic}{特性方程式}や総和型も同じ仕組みで扱える.

\subsection{母関数と数列の操作}
一般の\bookterm{sequence}{数列}$a_0,a_1,\ldots$では
\[
  A(x)=a_0+a_1x+a_2x^2+a_3x^3+\cdots
  =\sum_{n=0}^{\infty}a_nx^n
\]
とする.収束を仮定せず,各次数の係数を記録する形式的な記号として扱う.項の操作を代数の操作へ移すことが目的である.

たとえば $A(x)$ に $x$ を掛けると,
\[
  xA(x)=a_0x+a_1x^2+a_2x^3+\cdots
\]
となる.
最初に$0$を置き,項を一つずらした形になる.総和型では
$S_n=\sum_{k=0}^n a_k$とし,その\bookterm{generating}{母関数}を$S(x)$と書くと
\[
  S(x)=\frac{A(x)}{1-x}.
\]
これは,形式的な等比級数$\frac{1}{1-x}=1+x+x^2+\cdots$を掛けたとき,$x^n$の係数が$a_0+\cdots+a_n$になるからである.
逆に$(1-x)S(x)=A(x)$では,係数が$S_n-S_{n-1}=a_n$になる.\footnote{ただし$n=0$では$S_{-1}=0$として読む.}
総和型の「ずらして引く」は$1-x$倍になった.$a_1$からの和なら$a_0=0$として使える.

また,二つの\bookterm{sequence}{数列}の項を
\[
  c_n=\sum_{k=0}^{n}a_kb_{n-k}
\]
のように組み合わせる操作は\bookterm{convolution}{畳み込み}と呼ばれ,\bookterm{generating}{母関数}を掛けたときの $x^n$ の係数になる.
例えば全体を二つに分けて数える場合,それぞれの個数を掛けて分け方ごとに足す.このような和を,\bookterm{generating}{母関数}の積で表せる.

\begin{bookpoint}
  \textbf{POINT}\par
  各項を係数として母関数に記録すると,添字のずれは$x$倍,総和は$\frac{1}{1-x}$倍,畳み込みは積になる.
  ここでは級数の収束を仮定せず,各次数の係数を比較する.
\end{bookpoint}

\subsection{分母に現れる漸化式}
この翻訳をフィボナッチ\bookterm{sequence}{数列}で確かめてみよう.
\[
  F_0=0,\quad F_1=1,\qquad F_{n+2}=F_{n+1}+F_n
\]
とし,$F(x)=\sum_{n=0}^{\infty}F_nx^n$ と置く.
左辺の初めの項を取り出すと,
\[
  F(x)-x=\sum_{n=0}^{\infty}F_{n+2}x^{n+2}
\]
である.
\bookterm{recurrence}{漸化式} $F_{n+2}=F_{n+1}+F_n$ を代入すると,右辺の二つの和はそれぞれ $xF(x)$ と $x^2F(x)$ になる.
したがって
\[
  F(x)-x=xF(x)+x^2F(x)
\]
であり,
\[
  F(x)=\frac{x}{1-x-x^2}
\]
を得る.
\bookterm{recurrence}{漸化式}が\bookterm{generating}{母関数}の分母になった.

分母の方程式$1-x-x^2=0$に$r=\frac{1}{x}$を入れると,$r^2-r-1=0$になる.
分母が0になる点は特性方程式の解の逆数だった.行列の\bookterm{eigenvalue}{固有値}が1回の変化で何倍になるかを表すのに対し,母関数の分母は全項の係数関係を表す.

本編の符号に合わせ,$a_{n+2}+pa_{n+1}+qa_n=0$ $(n\ge0)$を考えても同じである.
$A(x)$の係数を比べると
\[
  (1+px+qx^2)A(x)=a_0+(a_1+pa_0)x.
\]
次数$2$以上の係数が\bookterm{recurrence}{漸化式}によってすべて0になるので,初期値の部分だけが右辺に残る.
\bookterm{characteristic}{特性方程式}$\lambda^2+p\lambda+q=0$の根を$\lambda_1,\lambda_2$とすると
\[
  1+px+qx^2=(1-\lambda_1x)(1-\lambda_2x).
\]
異なる0でない根の場合,\bookterm{partial}{部分分数分解}して
\[
  A(x)=\frac{C}{1-\lambda_1x}
  +\frac{D}{1-\lambda_2x}
\]
と書ける.
$\frac{1}{1-\lambda x}=\sum_{n=0}^\infty\lambda^n x^n$の係数を比べれば,$a_n=C\lambda_1^n+D\lambda_2^n$となり,本編の\bookterm{general}{一般項}と同じ形が得られる.
根の一つが0の場合も扱えるが,そのときは多項式の部分を分けて係数を読めばよい.

重解$r\ne0$の場合には分母が$(1-rx)^2$になる.
二つの等比級数を掛けると,$x^n$の係数は
\[
  \sum_{k=0}^n r^k r^{n-k}=(n+1)r^n
\]
だから
\[
  \frac1{(1-rx)^2}
  =\sum_{n=0}^\infty(n+1)r^n x^n.
\]
重解で$n$が現れた理由は,同じ因子の二乗からも読める.

\begin{bookpoint}
  \textbf{POINT}\par
  定数係数の同次線形漸化式は母関数の分母に,初期値は分子に現れる.
  分母の因子に応じて部分分数分解し,係数を読むと一般項に戻せる.重なった因子からは,累乗に掛かる$n$の式が現れる.
\end{bookpoint}

\BookImport{f-73b81d4e9a02}

\subsection{長期的な増え方}
\sectionref*{節：隣接3項間の漸化式を解く}のフィボナッチ数列の式$F_n=\frac{\varphi^n-\psi^n}{\sqrt5}$では,$|\psi|<1<\varphi$なので
\[
  F_n\sim\frac{\varphi^n}{\sqrt5}
\]
となる.
ここで $\sim$ は,二つの量の比が $1$ に近づくという意味である.
大きい絶対値の根が増大を支配する.同じ根を分母から読むことで,\bookterm{generating}{母関数}からも長期的な増大が読み取れる.

\subsection{畳み込みと積}
次は,項どうしを掛けて足す\bookterm{convolution}{畳み込み}を,数え上げの問題で使う.
因子の並びを変えず,二つずつ掛ける括弧の付け方を数えてみよう.
三つの因子$a,b,c$なら,$(ab)c$と$a(bc)$の二通りがある.
$n+1$個の因子に対する個数を$C_n$と書く.
一因子なら掛け算がなく$C_0=1$,二因子なら$C_1=1$,三因子なら$C_2=2$となる.
この列をカタラン数という.各掛け算を左右の分岐で描けば,左右を区別する二分木も数えている.

\paragraphbf{最後の掛け算で分ける}
$n+2$個の因子を掛け終える直前には,左側の積と右側の積が一つずつできている.
左が$k+1$個の因子なら,右は$n-k+1$個である.
左右それぞれの括弧の付け方を選ぶと,$C_kC_{n-k}$通りになる.
$k$は$0$から$n$まで動き,この分け方は最後の掛け算から一意に決まる.
従って,
\[
  C_0=1,\qquad C_{n+1}=\sum_{k=0}^{n}C_kC_{n-k}
\]
となる.左右を選ぶ積を全分け方で足すので,畳み込みになる.
\bookterm{generating}{母関数}$C(x)=\sum_{n=0}^{\infty}C_nx^n$では
\[
  C(x)=1+xC(x)^2
\]
という二次方程式になる.これを解くと
\[
  C(x)=\frac{1-\sqrt{1-4x}}{2x}
\]
である.
ここで平方根は,定数項が$1$となる形式的な級数として読む.\footnote{$H(x)=1+h_1x+h_2x^2+\cdots$とし,$H(x)^2=1-4x$の係数を合わせる.$x^n$の係数には$2h_n$と,先に決まった係数だけが現れるので,$h_n$を順に一意に求められる.ここでは数値を代入した級数の収束を使っていない.}
逆の符号では$x^{-1}$の項が生じるため,$C_0=1$から始まる母関数には合わない.
$x$の次数ごとに係数を求めると,
\[
  C(x)=1+x+2x^2+5x^3+\cdots
\]
となり,三因子の二通り,四因子の五通りも係数に戻る.

\paragraphbf{係数を取り出す}
では,$C_n$を$n$の式で表せないだろうか.
平方根を展開するのは高校範囲を超えるので,ここでは別の道を通ろう.
この計算では形式的な級数の微分を使うので,初めて読むときは飛ばしてよい.
脚注の$H(x)=\sqrt{1-4x}=\sum_{n=0}^{\infty}h_nx^n$を使う.$h_0=1,\ h_1=-2$である.
$H(x)^2=1-4x$の両辺を項ごとに微分すると$2H(x)H'(x)=-4$である.
両辺に$H(x)$を掛けて$H(x)^2=1-4x$を使えば,
\[
  (1-4x)H'(x)=-2H(x)
\]
が成り立つ.\footnote{級数を多項式と同じように項ごとに微分している.係数の並びとして扱う形式的な級数の立場では,収束を考えずにこの計算を正当化できる.}
両辺の$x^n$の係数を比べると,
\begin{align*}
  &(n+1)h_{n+1}-4nh_n=-2h_n\\
  &\therefore\quad h_{n+1}=\frac{4n-2}{n+1}h_n
\end{align*}
となる.なんと\sectionref*{節：階比型}である.
$C(x)=\frac{1-H(x)}{2x}$より$C_n=-\frac{h_{n+1}}{2}$なので,$C_n$も
\[
  C_n=\frac{2(2n-1)}{n+1}C_{n-1}\qquad(n\ge1)
\]
という階比型を満たす.書き下していけば
\begin{align*}
  C_n&=\prod_{k=1}^{n}\frac{2(2k-1)}{k+1}\\
  &=\frac{2^n\cdot1\cdot3\cdots(2n-1)}{(n+1)!}\\
  &=\frac{(2n)!}{n!\,(n+1)!}
\end{align*}
すなわち
\[
  C_n=\frac{1}{n+1}\binom{2n}{n}
\]
を得る.\footnote{最後の変形では$(2n)!=\{1\cdot3\cdots(2n-1)\}\cdot\{2\cdot4\cdots(2n)\}=1\cdot3\cdots(2n-1)\cdot2^nn!$を使った.}
総和を含む非線形の漸化式が,母関数を経由すると,本編の最初に学んだ階比型に帰着したのだ.
母関数は,扱える漸化式の範囲を大きく広げてくれる道具なのである.

\bookterm{nonlinear}{非線形}でも,和と積が母関数の操作に合えば代数方程式へ移せる.すべての非線形の漸化式に使えるわけではない.

この発想は,\bookterm{chebyshev}{チェビシェフ多項式}にも戻ってくる.
その\bookterm{recurrence}{漸化式}に\bookterm{generating}{母関数}を作ると,
\[
  \sum_{n=0}^{\infty}T_n(x)t^n
  =\frac{1-xt}{1-2xt+t^2}
\]
となる.
実際,$G(t)=\sum_{n=0}^\infty T_n(x)t^n$とし,$x$を定数として扱えば
\[
  (1-2xt+t^2)G(t)=1-xt.
\]
右辺は$T_0(x)=1$, $T_1(x)=x$に由来し,二次以上の係数は\bookterm{recurrence}{漸化式}によって打ち消される.
余弦の整数倍を表す列が,一つの有理関数にもまとまった.\bookterm{recurrence}{漸化式},\bookterm{trigonometric}{三角関数},\bookterm{generating}{母関数}は同じ列の三つの表現である.

\paragraphbf{別の広がり：階乗の番号を広げる}
この部分では広義積分など数学IIIの内容を使う.初めて読むときは飛ばしてもよい.
母関数は各項を係数に記録した.次は,項の\textbf{番号そのもの}を整数以外へ広げる.
階乗は$u_0=1$, $u_{n+1}=(n+1)u_n$という\bookterm{ratio}{階比}型の例である（\booknavref{解法：階比型：f(n)が1次式の積の場合}{1次式の積の場合}）.
その関係を広げるガンマ関数は,正の$x$に対して
\[
  \Gamma(x)=\int_0^\infty t^{x-1}e^{-t}\,dt
\]
と定められる.
部分積分をすると
\[
  \Gamma(x+1)=x\Gamma(x)
\]
が現れる.最初の値も
\[
  \Gamma(1)=\int_0^\infty e^{-t}\,dt=1
\]
なので,整数$n\ge0$では$\Gamma(n+1)=n!$となる.
番号を広げた階乗が積分や確率につながる.\bookterm{recurrence}{漸化式}から関数を作る別の例である.

\subsection{確率と整数への応用}
確率に目を向けると, 状態の移り変わりも\bookterm{recurrence}{漸化式}で記録できる.
晴れ・雨の確率を並べた\bookterm{column-vector}{列ベクトル}を$\boldsymbol p_n$とする.晴れ,雨の翌日に晴れる確率がそれぞれ$0.8,0.3$なら
\[
  \boldsymbol p_{n+1}
  =
  \begin{pmatrix}0.8&0.3\\0.2&0.7\end{pmatrix}
  \boldsymbol p_n
  =P\boldsymbol p_n
\]
と書ける.
各列の和は$1$なので,確率の合計は$1$に保たれる.
晴れの確率を$u_n$とすれば,雨の確率は$1-u_n$だから
\[
  u_{n+1}=0.8u_n+0.3(1-u_n)=0.5u_n+0.3.
\]
\sectionref*{節：特性方程式型}で扱った定数項をもつ式そのもので,\bookterm{fixed}{不動点}は$0.6$である.
したがって
\[
  u_n-0.6=(u_0-0.6)(0.5)^n
\]
となり,晴れ・雨の確率は$(0.6,0.4)$へ近づく.
\bookterm{matrix}{行列}では,$(1,1)P=(1,1)$が確率の合計を保つことを表し,\bookterm{eigenvalue}{固有値}$0.5$は\bookterm{fixed}{不動点}からのずれが1回ごとに何倍になるかを表す.
不動点と等比型が,確率の長期予測にも使えた.

数論では, ユークリッドの互除法が余りを次々に作る漸化式として表せる.
二つの正の整数$a_0,a_1$から始める.
$a_{n-1}$ を $a_n$ で割った商を $q_n$, 余りを $a_{n+1}$ とすれば,
\[
  a_{n-1}=q_na_n+a_{n+1},
  \qquad 0\le a_{n+1}<a_n
\]
となる.
余りが0なら計算を止め,そうでなければ2つの数を$(a_n,a_{n+1})$へ置き換える.
$a_{n+1}=a_{n-1}-q_na_n$と逆向きの式から,共通の約数は変わらない.従って
\[
  \gcd(a_{n-1},a_n)=\gcd(a_n,a_{n+1}).
\]
正の余りは小さくなるので有限回で止まる.
余り$a_{n+1}$が$0$になったときの除数$a_n$,すなわち最後の0でない項が最大公約数となる.
確率の合計と同じく,変わらない量から計算の意味を説明できる.

整数係数・整数初期値の隣接3項型では,各項を整数$m\ge2$で割った余りだけを残してみよう.
状態 $(a_n,a_{n+1})$ の各成分は $0,1,\ldots,m-1$ のいずれかなので, 状態は高々 $m^2$ 通りしかない.
次の状態が今の状態だけで決まるなら, いつか同じ状態が現れ, その後は同じ並びを繰り返す.
したがって,余りの列はある項以降で周期的になる.
たとえばフィボナッチ\bookterm{sequence}{数列}を $2$ で割った余りは
\[
  0,1,1,0,1,1,0,1,1,\ldots
\]
と周期 $3$ で最初から繰り返す.
一方,漸化式を逆向きにたどれない場合には,最初から周期的になるとは限らない.
たとえば $a_1=a_2=1$,$a_{n+2}=2a_{n+1}$ を $4$ で割った余りは
\[
  1,1,2,0,0,0,\ldots
\]
となり,$a_4$ 以降は周期 $1$ で繰り返すが,初めの $1,1,2$ は繰り返されない.
同じ漸化式でも,実数の増大と余りの周期では見える性質が変わる.

本編の\bookterm{trigonometric}{三角関数}型に似た二次式でも,同じ解法が使えるとは限らない.それでも増大は評価できる.
たとえば
\[
  a_0=0,\qquad a_{n+1}=a_n^2+1
\]
では $0,1,2,5,26,\ldots$ と増える.
$n\ge2$ では $a_n\ge2$ であり, $a_{n+1}\ge a_n^2$ だから, 帰納的に
\[
  a_n\ge 2^{\,2^{n-2}}
\]
となる.
一般項を求めずに,急速な増大を評価できた.
また, \bookterm{sequence}{数列}が単調で有界と分かれば収束し, 漸化式が $a_{n+1}=f(a_n)$ で $f$ が連続なら, \bookterm{limit}{極限} $L$ は $L=f(L)$ を満たす.
単調性, 有界性, \bookterm{fixed}{不動点}, 増加速度は, \bookterm{general}{一般項}とは別の問いに答えてくれる.

行列・差分・母関数は,本編の操作を別の表現へ移す道具だった.
不動点からのずれ,\bookterm{left-eigen}{左固有ベクトル}で取り出す量,添字のずれは,それぞれ同じ漸化式の別の表し方である.
一般項に加え,増大・確率・周期・保存される量も,見たい性質に合わせて調べられる.

\begin{bookpoint}
  \textbf{POINT}\par
  調べたい性質に合わせ,数列全体を母関数で表したり,確率や余りなどの量に注目したりする.
  極限の候補を不動点の式から得た場合も,収束そのものは単調性・有界性などで別に確かめる.
\end{bookpoint}


  % BOOK-FILE-ID: f-45a181b08d92
\section{演習を別の視点から}\label{節：演習を別の視点から}
\bookproblemref{practice-entrance}{9}{発展問9},
\bookproblemref{practice-entrance}{4}{問4},
\bookproblemref{practice-entrance}{3}{問3}で使った操作を,不変量・回転・母関数という別の視点から見直す.
演習で解いた問題の裏にも,本編と同じ数学が隠れていたことが分かるだろう.

\subsection{整数方程式と不変量}
発展問9では,$u_0=3$, $u_1=3$, $u_{n+2}=3u_{n+1}-u_n$から,
$(3,u_n,u_{n+1})$が$x^2+y^2+z^2=xyz$の解になることを示した.
$n=0$では
\[
  u_n^2+u_{n+1}^2+9=3u_nu_{n+1}
\]
が成り立つ.一般の$x,y$に対して
\[
  y^2+(3y-x)^2-3y(3y-x)=x^2+y^2-3xy
\]
だから,$u_n^2+u_{n+1}^2-3u_nu_{n+1}$は漸化式を用いる前後で変わらず$-9$に保たれる.
$u_{n+2}=3u_{n+1}-u_n$を使えば,$u_nu_{n+2}=3u_nu_{n+1}-u_n^2=u_{n+1}^2+9$となり,
\[
  u_{n+2}=\frac{u_{n+1}^2+9}{u_n}
\]
とも書ける.正値性は$u_0=u_1=3$と,
$u_2=6$以降の増加から確認できる.
\sectionref*{節：変わらない量を使う}の定数$1$を$9$へ変えた漸化式だった.
元の問題は\bookterm{linear}{線形}の漸化式から保存を示し,本編の例は保存から線形の漸化式を取り出した.

\begin{bookpoint}
  \textbf{POINT}\par
  隣り合う項の組に対して,更新しても変わらない式を探し,初期値でその値を決める.
  保存される関係を使うと,非線形の式から線形の漸化式を取り出せる場合がある.
\end{bookpoint}

\subsection{中点・回転・面積}
発展問4では,平行四辺形の各辺上で同じ比率$t$ $(0<t<1)$の分点を取り,
新しい平行四辺形を作った.連続する辺ベクトルを$u,v$とする.
頂点についての式を引き算すれば,
\[
  u'=(1-t)u+tv,\qquad v'=-tu+(1-t)v
\]
である.

\paragraphbf{正方形の中点を結ぶ操作}
正方形で,向きを反時計回りに取った辺を複素数として$v=iu$と書けるときは,
\[
  u'=((1-t)+it)u
\]
である.$t=\frac12$なら拡大率は$\frac1{\sqrt2}$,回転角は$\frac\pi4$になる.
これが問題で中点を結んだときの辺長と向きに一致する.

\par\medskip\noindent\textbf{別解：平行四辺形を行列で見る}\par\nobreak\smallskip
正方形という条件を外し,最初に求めた$u',v'$の式へ戻る.
辺の組を混ぜる\bookterm{matrix}{行列}は
\[
  M=\begin{pmatrix}1-t&t\\-t&1-t\end{pmatrix}
\]
となる.\bookterm{matrix}{行列}式は$\rho=(1-t)^2+t^2$である.
$u=(u_1,u_2)$, $v=(v_1,v_2)$と書くと,底辺と高さから平行四辺形の面積は$|u_1v_2-u_2v_1|$となる.$u',v'$の式を代入すると
\BookMathLayout{\[
    u'_1v'_2-u'_2v'_1
    =\bigl((1-t)^2+t^2\bigr)(u_1v_2-u_2v_1).
  \]}{\[
    \begin{aligned}
      u'_1v'_2-u'_2v'_1\\
      {}=\bigl((1-t)^2+t^2\bigr)(u_1v_2-u_2v_1).
    \end{aligned}
  \]}
$\rho>0$なので,絶対値を取った面積も一回ごとに$\rho$倍となる.
この行列式の値は,\bookterm{eigenvalue}{固有値}$1-t\pm it$の絶対値の二乗でもある.

\begin{bookpoint}
  \textbf{POINT}\quad \bookterm{matrix}{行列}が作用する対象\par
  辺の組を混ぜる\bookterm{matrix}{行列}を,平面上の相似変換(長さを一様に拡大・縮小し,角度を保つ変換)と読むには条件が要る.
  面積が何倍になるかは一般の平行四辺形でも分かるが,上の回転の説明には直交・等長の辺を使った.
\end{bookpoint}

\subsection{一定になる畳み込み和}
発展問3の$c_0=1$, $c_{n+1}=\frac{2n+1}{2n+2}c_n$から,
\[
  c_n=\frac1{4^n}\binom{2n}{n}
\]
を得た.その\bookterm{generating}{母関数}を$C(x)=\sum_{n\ge0}c_nx^n$とする.
$C(x)^2$の$x^n$の係数が,求める和$\sum_{k=0}^n c_kc_{n-k}$である.
ここでは積の微分法則と,高校範囲を超える形式的な母関数の微分を使う.初めて読むときは飛ばしてもよい.
形式的な微分は$x^n\mapsto nx^{n-1}$と定め,関数としての収束は仮定しない.
元の\bookterm{recurrence}{漸化式}を$(n+1)c_{n+1}=(n+\frac12)c_n$と書けば,
\[
  (1-x)C'(x)=\frac12C(x)
\]
を得る.積の係数は有限和なので,係数を比較することで通常の積の微分法則も使える.
$D(x)=C(x)^2=\sum_{n\ge0}d_nx^n$と置くと,
\[
  (1-x)D'(x)=2C(x)(1-x)C'(x)=D(x).
\]
$x^n$の係数を比較すれば$(n+1)d_{n+1}-nd_n=d_n$,すなわち$d_{n+1}=d_n$である.
$d_0=c_0^2=1$なので,
\[
  C(x)^2=\frac1{1-x}=\sum_{n\ge0}x^n
\]
の$x^n$の係数から
\[
  \sum_{k=0}^n c_kc_{n-k}=1
\]
を得る.有限和の変形と母関数の積は,同じ\bookterm{convolution}{畳み込み}の二つの扱い方である.

\begin{bookpoint}
  \textbf{POINT}\quad 別解の使い分け\par
  別の表現は,和・面積・不変量の意味を見せる.元の条件と答えへ戻って対応を確かめる.
\end{bookpoint}


\end{bookcolumns}
